% ============================================================
%  content.tex --- Informe 01: De lo Actuado
%  Proyecto: Dise\~no El\'ectrico Conjunto Residencial 65 Casas
%  Asignatura: Instalaciones El\'ectricas
%  Equipo: Jerson Gamboa, Daniel Fernando, Jhon
% ============================================================

% ─────────────────────────────────────────────────────────────
\section{Introducci\'on}
% ─────────────────────────────────────────────────────────────

El presente informe describe la metodolog\'ia de trabajo adoptada
por el equipo durante el desarrollo del proyecto de dise\~no de
la infraestructura el\'ectrica de un conjunto residencial cerrado
de 65 viviendas unifamiliares en Bucaramanga, Santander. El documento
da cuenta de la organizaci\'on interna del grupo, la distribuci\'on
de responsabilidades, las herramientas utilizadas y los mecanismos
de control de calidad implementados para garantizar la consistencia
t\'ecnica entre los entregables del proyecto.

El proyecto abarca el dise\~no completo de la red de distribuci\'on
en media tensi\'on a $13{,}2\text{ kV}$, la subred de baja tensi\'on
en $127/220\text{ V}$, el sistema de puesta a tierra, las memorias
de c\'alculo, las plantillas de verificaci\'on normativa y los planos
constructivos en AutoCAD, todo ello enmarcado en el cumplimiento del
Reglamento T\'ecnico de Instalaciones El\'ectricas
(RETIE) y la Norma T\'ecnica Colombiana NTC~2050.

% ─────────────────────────────────────────────────────────────
\section{Descripci\'on General del Proyecto}
% ─────────────────────────────────────────────────────────────

El conjunto residencial objeto del dise\~no cuenta con las siguientes
caracter\'isticas t\'ecnicas generales:

\begin{tabular}{ll}
  \toprule
  \bfseries Par\'ametro & \bfseries Valor \\
  \midrule
  N\'umero de viviendas         & 65 unidades unifamiliares \\
  Tensi\'on de suministro MT    & $13{,}2\text{ kV}$ \\
  Tensi\'on de distribuci\'on BT & $127/220\text{ V}$ (trif\'asico) \\
  Tipo de red                  & 100\% subterr\'anea \\
  Conductor utilizado          & THHN en canaleta subterr\'anea \\
  Topolog\'ia                  & Radial con selectividad por caja \\
  Medidores por vivienda       & 2 (residencial + veh\'iculo el\'ectrico) \\
  Medici\'on en frontera       & Macromedidor general en MT \\
  \bottomrule
\end{tabular}

\vspace{6pt}

La elecci\'on de la topolog\'ia radial, en lugar del esquema
en anillo, respondi\'o a restricciones regulatorias del operador
de red, y ofrece como ventaja operativa la selectividad ante fallas:
un evento el\'ectrico en cualquier vivienda o sector queda aislado
en su caja de protecci\'on correspondiente sin afectar al resto
del conjunto \citep{icontec2050}.

% ─────────────────────────────────────────────────────────────
\section{Organizaci\'on del Equipo de Trabajo}
% ─────────────────────────────────────────────────────────────

El equipo estuvo conformado por tres estudiantes de la asignatura
Instalaciones El\'ectricas, quienes distribuyeron las
responsabilidades del proyecto seg\'un sus fortalezas individuales,
manteniendo en todo momento un flujo de trabajo colaborativo y
con revisi\'on cruzada de los entregables.

\subsection{Roles y Responsabilidades Principales}

\begin{tabular}{lp{4cm}p{7cm}}
  \toprule
  \bfseries Integrante &
  \bfseries Rol principal &
  \bfseries Actividades lideradas \\
  \midrule
  Jerson Gamboa &
    Dise\~no el\'ectrico y normativa &
    Lider\'o el desarrollo de la malla de tierra,
    la verificaci\'on de regulaciones RETIE y
    la redacci\'on de las memorias de c\'alculo \\[6pt]
  Daniel Fernando &
    Gesti\'on documental y estructuraci\'on &
    Lider\'o la organizaci\'on inicial de carpetas,
    la estructuraci\'on de archivos del proyecto
    y la coordinaci\'on de entregables \\[6pt]
  Jhon &
    Dise\~no el\'ectrico y normativa &
    Apoy\'o junto a Jerson el desarrollo de
    regulaciones, memorias t\'ecnicas y
    diligenciamiento de plantillas de verificaci\'on \\
  \bottomrule
\end{tabular}

\vspace{6pt}

Cabe destacar que, si bien cada integrante tuvo un \'area de liderazgo
definida, la naturaleza del proyecto demand\'o participaci\'on activa
de todo el equipo en las actividades transversales. En particular,
el diligenciamiento de las plantillas de c\'alculo en Excel y el
dibujo de los planos en AutoCAD fueron actividades compartidas y
distribuidas de manera equitativa entre los tres integrantes.

% ─────────────────────────────────────────────────────────────
\section{Metodolog\'ia de Trabajo}
% ─────────────────────────────────────────────────────────────

\subsection{Fase 1: Organizaci\'on y Estructuraci\'on del Proyecto}

La primera etapa del proceso consisti\'o en definir la arquitectura
documental del proyecto. Daniel Fernando lider\'o esta fase,
estableciendo la estructura de carpetas, la nomenclatura de archivos
y el flujo de versiones de los documentos. Esta decisi\'on temprana
result\'o determinante para mantener la trazabilidad de los
entregables a lo largo de todo el desarrollo y evitar la
duplicaci\'on o p\'erdida de informaci\'on.

Se defini\'o la siguiente estructura de entregables:

\begin{enumerate}
  \item Memorias de c\'alculo (dimensionamiento de conductores,
    protecci\'on y ca\'idas de tensi\'on)
  \item Plantillas de verificaci\'on normativa en Excel
  \item Planos constructivos en AutoCAD
  \item Informe de lo actuado (presente documento)
  \item Descripci\'on general para inversionistas
  \item Informe de cumplimiento RETIE
\end{enumerate}

\subsection{Fase 2: Desarrollo T\'ecnico}

Con la estructura documental definida, el equipo procedi\'o al
desarrollo t\'ecnico del proyecto. Jerson Gamboa y Jhon lideraron
esta fase, con enfoque en tres frentes simult\'aneos:

\textbf{Malla de tierra.} Se dise\~n\'o el sistema de puesta a
tierra del conjunto siguiendo los criterios de la NTC~2050 y el
RETIE, garantizando valores de resistencia dentro de los l\'imites
exigidos para instalaciones de media tensi\'on
\citep{minminas2013retie}.

\textbf{C\'alculo de regulaciones y memorias.} Se elaboraron las
memorias de c\'alculo para el dimensionamiento de conductores THHN,
la verificaci\'on de ca\'idas de tensi\'on admisibles y la
selecci\'on de elementos de protecci\'on (fusibles y
interruptores) para cada nivel de la red.

\textbf{Plantillas de verificaci\'on en Excel.} Todo el equipo
particip\'o en el diligenciamiento de las plantillas de c\'alculo,
que permit\'ian cruzar autom\'aticamente los valores obtenidos con
los l\'imites normativos. Este trabajo fue desarrollado de forma
paralela entre los tres integrantes para acelerar los tiempos
de entrega.

\subsection{Fase 3: Dise\~no Gr\'afico (Planos AutoCAD)}

La elaboraci\'on de los planos constructivos se distribuy\'o
equitativamente entre los tres integrantes del equipo. Se
trabajaron los siguientes tipos de plano:

\begin{itemize}
  \item Plano de distribuci\'on general de la red MT
  \item Plano de subred de baja tensi\'on por manzana
  \item Diagrama unifilar general del conjunto
  \item Plano de malla de puesta a tierra
  \item Esquema del sistema de medici\'on
    (macromedidor MT + submedidores BT)
\end{itemize}

La distribuci\'on equitativa del trabajo gr\'afico garantiz\'o
que todos los integrantes tuvieran dominio del modelo completo
del proyecto y pudieran responder por cualquier elemento del
dise\~no ante una evaluaci\'on o revisi\'on externa.

\subsection{Fase 4: Control de Calidad --- Revisi\'on Cruzada}

Antes de la entrega formal de cada entregable, el equipo
implement\'o una metodolog\'ia de \textbf{revisi\'on cruzada}:
cada documento o plano elaborado por un integrante fue revisado
por otro miembro del equipo diferente al autor original.

Este mecanismo permiti\'o detectar inconsistencias entre las
memorias de c\'alculo y los planos, errores en el diligenciamiento
de plantillas y omisiones en la verificaci\'on normativa, antes
de que los documentos fueran consolidados como versi\'on final.
La revisi\'on cruzada es una pr\'actica est\'andar en proyectos
de ingenier\'ia de alta responsabilidad y fue adoptada conscientemente
por el equipo como garant\'ia de calidad del producto entregado
\citep{icontec2050}.

% ─────────────────────────────────────────────────────────────
\section{Herramientas Utilizadas}
% ─────────────────────────────────────────────────────────────

\begin{tabular}{lll}
  \toprule
  \bfseries Herramienta & \bfseries Uso & \bfseries Fase \\
  \midrule
  Microsoft Excel &
    Plantillas de c\'alculo y verificaci\'on normativa &
    2 y 4 \\
  AutoCAD &
    Elaboraci\'on de planos constructivos &
    3 \\
  \LaTeX{} / Tectonic &
    Redacci\'on de informes en formato APA &
    1, 2, 3 y 4 \\
  Normativa RETIE &
    Marco regulatorio de referencia para el dise\~no &
    2 y 4 \\
  NTC 2050 &
    C\'odigo el\'ectrico nacional de referencia &
    2 y 4 \\
  \bottomrule
\end{tabular}

% ─────────────────────────────────────────────────────────────
\section{Resultados del Proceso}
% ─────────────────────────────────────────────────────────────

Como resultado del proceso descrito, el equipo entreg\'o un
conjunto completo y consistente de documentos t\'ecnicos que
cubren la totalidad del dise\~no el\'ectrico del proyecto.
Los entregables finales son coherentes entre s\'i: los valores
calculados en las memorias se reflejan fielmente en los planos,
y las plantillas de Excel acreditan el cumplimiento normativo
de cada elemento dise\~nado.

El elemento de mayor valor agregado del proceso fue la
implementaci\'on del sistema de \textbf{macromedici\'on en
media tensi\'on con balance de energ\'ia}, que permite calcular
el consumo de las zonas comunes del conjunto sin necesidad de
un medidor f\'isico adicional, mediante la expresi\'on:

\begin{equation}
  E_{\text{zonas comunes}} =
    E_{\text{frontera MT}} -
    \sum_{i=1}^{130} E_{\text{medidor}_i}
\end{equation}

Este esquema, que debe ser radicado formalmente ante la
Electrificadora de Santander (ESSA) en la etapa de
legalizaci\'on, representa una soluci\'on t\'ecnica robusta,
transparente y auditables para la administraci\'on del
conjunto residencial \citep{essa2020}.

% ─────────────────────────────────────────────────────────────
\section{Conclusiones}
% ─────────────────────────────────────────────────────────────

El desarrollo del proyecto permiti\'o al equipo integrar los
conocimientos te\'oricos de la asignatura Instalaciones
El\'ectricas con la pr\'actica real del dise\~no de redes de
distribuci\'on en media y baja tensi\'on.

La metodolog\'ia de trabajo adoptada --- con roles definidos,
participaci\'on transversal en actividades comunes y revisi\'on
cruzada de entregables --- demostr\'o ser eficiente para
garantizar la calidad t\'ecnica del producto final en un
contexto de trabajo en equipo acad\'emico.

El equipo concluye que la correcta estructuraci\'on documental
desde la fase inicial del proyecto es un factor cr\'itico
de \'exito, ya que facilita la trazabilidad de decisiones
de dise\~no, reduce los tiempos de correcci\'on y asegura
la consistencia entre memorias, plantillas y planos
constructivos.